\documentclass[11pt]{article}
\usepackage[a4paper,margin=24mm]{geometry}
\usepackage{fontspec,amsmath,amssymb,amsthm,fvextra,hyperref,longtable}
\setmainfont{lmroman10-regular.otf}[BoldFont=lmroman10-bold.otf,ItalicFont=lmroman10-italic.otf]
\setmonofont{JuliaMono-Regular.ttf}[Path=fonts/,Scale=0.90]
\hypersetup{colorlinks=true,linkcolor=blue,urlcolor=blue}
\newtheorem{theorem}{Theorem}
\newtheorem{lemma}[theorem]{Lemma}
\newcommand{\N}{\mathbb N}
\newcommand{\R}{\mathbb R}
\title{Collatz target density: statement, proof chain, and Lean verification}
\author{Local formalization report}
\date{12 September 2026}
\begin{document}
\maketitle
\begin{abstract}
This report documents the target-density theorem in the accompanying Lean
project. For a positive integer target, its Collatz predecessors have positive
lower natural density if and only if the target is not divisible by three.
The argument combines elementary inverse arithmetic and nonperiodic seed
selection with imported analytic density machinery. The mathematical
exposition explicitly identifies that substantial imported input; it is not
a new self-contained derivation of all its analytic estimates. Complete local
Lean modules in the dependency chain are reproduced at the end. This result
does not prove that every positive integer reaches one, and no claim of
historical priority is made.
\end{abstract}
\tableofcontents
\section{The precise assertion}
Define the ordinary Collatz map on nonnegative integers by
\[
 T(n)=\begin{cases}n/2&n\text{ even},\\3n+1&n\text{ odd}.\end{cases}
\]
Iteration includes the zeroth iterate, $T^0(n)=n$. For $a>0$ put
\[
 R_a=\{n\in\N:n>0,\ \exists k\in\N\ T^k(n)=a\},\qquad
 P_a(X)=\#(R_a\cap[0,X)).
\]
Thus the cutoff is strict and the starting integer zero is excluded.
\begin{theorem}[Public target-density theorem]
For every positive integer $a$,
\[
 \bigl(\exists c>0\ \exists X_0\in\N\ \forall X\ge X_0:
 P_a(X)\ge cX\bigr)\quad\Longleftrightarrow\quad 3\nmid a.
\]
The constant and cutoff may depend on $a$.
\end{theorem}
Equivalently, the lower natural density $\liminf_{X\to\infty}P_a(X)/X$
is positive exactly for these targets. This equivalence with limit notation
is elementary; the literal Lean endpoint uses the quantified inequality.
For $a=1$ the theorem supplies a positive fraction of convergent starts,
not the assertion that every start converges.

\section{Arithmetic supply of inverse seeds}
For odd positive $n$, let
\[
 S(n)=\frac{3n+1}{2^{v_2(3n+1)}}
\]
be the Syracuse map. Every Syracuse step is a finite sequence of ordinary
Collatz steps. Write $\operatorname{odd}(m)=m/2^{v_2(m)}$ for $m>0$.

\begin{lemma}[One-step inverse arithmetic]
If $a$ is odd and $3\nmid a$, there is an odd positive $c$ with $S(c)=a$.
\end{lemma}
\begin{proof}
If $a\equiv1\pmod3$, take $c=(4a-1)/3$; if $a\equiv2\pmod3$, take
$c=(2a-1)/3$. Each quotient is a positive odd integer, and respectively
$3c+1=4a$ or $3c+1=2a$. Removing all powers of two gives $a$.
\end{proof}

Let $u_j=(4^j-1)/3$. For $j\ge1$ it is odd and $S(u_j)=1$.
The identities
\[
 u_{m+n}=u_m+4^m u_n,\qquad
 v_3(4^n-1)=1+v_3(n)\quad(n>0)
\]
imply $3^q\mid u_n$ if and only if $3^q\mid n$. Consequently
\[
 u_m\equiv u_{m+n}\pmod{3^q}\quad\Longleftrightarrow\quad 3^q\mid n.
\]
The map $j\mapsto u_j\pmod{3^q}$ permutes the $3^q$ residues as $j$
ranges from $0$ to $3^q-1$. Increasing $j$ by arbitrarily large multiples
of $3^q$ preserves its residue and makes $u_j$ arbitrarily large.
The valuation identity is the standard lifting-the-exponent calculation;
the Lean proof uses the multiplicity theorem in Mathlib. This residue
argument is copied, with attribution, from the upstream source.

Fix the child $c$ above and define
\[
 F_c(x)=(3c+1)x+c.
\]
Then
\[
 3F_c(x)+1=(3c+1)(3x+1),\qquad
 S(F_c(x))=\operatorname{odd}(3c+1)\operatorname{odd}(3x+1)=a
\]
whenever $S(x)=1$. Also $F_c(x)$ is odd, is larger than $x$, and its slope
$3c+1$ is coprime to $3^q$. Hence $F_c$ permutes residues modulo $3^q$.
For every $q$, residue $r$, and size bound $B$, there is therefore an odd
$M>B$ with $M\equiv r\pmod{3^q}$ and $S(M)=a$.

\section{Removing periodic seeds without assuming convergence}
Call $x$ nonperiodic for a function $f$ if $f^k(x)\ne x$ for every $k>0$.
This permits a point that eventually enters a cycle: only the point itself
must not lie on a cycle.
\begin{lemma}
If $x\ne y$ and $f(x)=f(y)$, at least one of $x,y$ is nonperiodic.
Furthermore, a forward orbit hits a nonperiodic point at most once.
\end{lemma}
\begin{proof}
If $x,y$ have respective positive periods $p,q$, both are fixed by $f^{pq}$.
Applying $f^{pq-1}$ to $f(x)=f(y)$ gives $x=y$, a contradiction.
If $f^d(z)=f^e(z)=x$ with $d<e$, then $f^{e-d}(x)=x$,
contradicting nonperiodicity.
\end{proof}
Choose two increasingly large children in any specified ternary residue
class. They are distinct and have the same image $a$, so one is
nonperiodic. This proves the exact arithmetic input
\[
 \begin{split}
 \operatorname{SeedSupply}(a):\quad
 &\forall q,B\in\N\ \forall r\in\mathbb Z/3^q\mathbb Z\ \exists M>B:\\
 &M\text{ odd},\quad S(M)=a,\quad M\equiv r\pmod{3^q},\quad
 M\text{ nonperiodic}.
 \end{split}
\]
No convergence of $a$ to one occurs in this argument.

\section{The substantial analytic input and its use}
The project imports Lech Mazur's Apache-2.0 positive-density source package,
identified in the local modules by corrected commit
\texttt{0aef8b0cfaaf6959500063472f2e328c60df8d54}.
Its analytic estimates are a substantive part of this proof, not an
elementary consequence of residue coverage. Residue coverage alone would
not justify positive natural density.

Here is the chain used in \texttt{Activation}, \texttt{SourceCharge}, and
\texttt{Assembly}. The terminology ``mass'' denotes the weighted inverse-path
quantities defined in the imported project; it does not denote a count of
distinct integers until the counting estimate is applied.

\begin{enumerate}
\item The imported finite-residue averaging estimate chooses a residue with
marked mass at least $2/3$ times a probability product. The imported product
has a positive lower bound $p$. Seed supply realizes that residue by an
arbitrarily large nonperiodic physical seed reaching $a$.
\item Imported variation bounds are uniform in the seed and have tails
tending to zero. Choose a generation with tail less than $p/3$ and then
choose the seed. The marked sequence for this one fixed seed consequently
has a positive limit. This order of choices is essential: a different seed
at each generation would not suffice.
\item The terminal-mark construction gives a fixed $\alpha>0$ and an eventual
lower bound $\alpha$ on core terminal unit mass. Nonperiodicity supplies
unique hitting times for the source-charge estimate. If $P>0$ is parent
source potential and $C>0$ the imported rate constant, choose an integer
$m>\max(1,88PC/\alpha)$. Then
\[
 44PC/m^2\le\alpha/2.
\]
With $H=(8/9)3^m$, the rate inequality gives eventual full terminal mass
at least $\eta=\alpha/(2H)>0$.
\item The imported interval synchronization theorem makes this estimate
available for every sufficiently large cutoff, not merely a subsequence.
For $x=Y/32$, the terminal sources lie in $[x,32x)$, are odd, and reach $a$.
The source-charge bound gives
\[
 x\eta\le x\,\mathrm{Mass}\le P\,\#\mathrm{Sources}
 \le P\,\#\{n<Y:n\text{ odd and }S^k(n)=a\text{ for some }k\}.
\]
Thus the last count is at least $\eta Y/(32P)$ for every sufficiently
large $Y$.
\end{enumerate}
The formal theorem encapsulating this chain is
\texttt{positive\_density\_of\_seed\_supply}. The source listings below
contain its local proof and all its local dependencies. Imported analytic
modules are retained in the accompanying vendor directory; their full
contents are not printed here. Accordingly this exposition is a proof
relative to explicitly identified, formally checked imported lemmas, not
a replacement for their complete development.

\section{Ordinary Collatz targets and the sharp exception}
For any positive $a$ with $3\nmid a$, including even $a$, the same formulas
produce an odd $c$ with $3c+1=2^v a$. If $3\mid c$, replace $c$ by
$c'=4c+1$. Then
\[
 3c'+1=4(3c+1)=2^{v+2}a,\qquad c'\equiv1\pmod3.
\]
We now have an odd child coprime to three which reaches $a$ under $T$.
Its odd Syracuse predecessor set has positive lower density by the preceding
argument and is contained in $R_a$. Monotonicity of counting proves the
forward implication of the theorem.

Suppose instead $3\mid a$. An odd Collatz step has image congruent to one
modulo three. Working backwards from a multiple of three, every step must
therefore be an even step, whose unique inverse is doubling. Induction
on the number of steps proves
\[
 T^k(n)=a\text{ for some }k\quad\Longleftrightarrow\quad
 n=2^j a\text{ for some }j\ge0.
\]
In particular $P_a(2^k a)=k$. A positive eventual lower bound $P_a(X)\ge cX$
would force $c2^k a\le k$ for all large $k$, which is impossible.
The Lean proof uses $k^2\le2^k$ for $k\ge4$ and chooses $kc>1$.
This proves the converse and completes the theorem.

\section{What this says about counterexamples}
Let $B=\{n>0:n\text{ never reaches }1\}$. Every positive orbit reaches
some integer coprime to three: remove initial factors of two and, if
necessary, perform one odd step. If $B$ is nonempty, choose such an $a$
on a bad orbit. Then $a$ is bad. Every predecessor of $a$ is bad as well:
an orbit that reaches one remains on $1,4,2$, so every later state also
reaches one. Therefore $R_a\subseteq B$ and
\[
 B\ne\varnothing\quad\Longrightarrow\quad
 \exists c>0\ \exists X_0\ \forall X\ge X_0:\ \#(B\cap[0,X))\ge cX.
\]
Consequently full Collatz is equivalent in this project to
\[
 \forall\varepsilon>0\ \forall X_0\ \exists X\ge X_0,
 \quad X>0,\quad \#(B\cap[0,X))<\varepsilon X.
\]
That sparsity condition has not been proved. Positive density of each
individual predecessor set does not imply that their union exhausts every
integer, or that a bad orbit cannot exist. The formal result is useful
structural information, but is not a solution of the Collatz conjecture.

\section{Verification and reproducibility}
The verification run completed successfully on 12 September 2026.
The toolchain was Lean 4.30.0-rc2. Run from the project directory:
\begin{Verbatim}[breaklines=true]
LEAN_NUM_THREADS=4 python3 scripts/verify.py
\end{Verbatim}
The script reported \texttt{verified}. It checked 1,485 upstream source
entries against their manifest hashes, checked nine dependency revisions
against the upstream lock file, recorded local source SHA-256 hashes,
and completed the build of both libraries (5,105 build jobs, including
replayed jobs). It then audited 255 listed theorem endpoints. The public
target-density theorem and counterexample dichotomy use only
\texttt{propext}, \texttt{Classical.choice}, and \texttt{Quot.sound}.
No audited endpoint used an axiom outside this standard set. The source
escape scan also passed; in particular the scanned proof sources contain
no executable \texttt{sorry} or additional axiom declaration.

The machine-readable record is \texttt{verification/local-report.json};
build and axiom logs are in the same directory. Every one of the 13 local
modules reproduced below was part of this build. The appendix contains
1,283 source lines and was compared byte-for-byte with those modules.
It relies on the project's pinned imported libraries, so compilation of
isolated snippets without those imports is not the verification procedure.

To regenerate this PDF from its output directory, run XeLaTeX twice:
\begin{Verbatim}[breaklines=true]
xelatex -interaction=nonstopmode -halt-on-error collatz_target_density_report.tex
xelatex -interaction=nonstopmode -halt-on-error collatz_target_density_report.tex
\end{Verbatim}
The adjacent \texttt{fonts/JuliaMono-Regular.ttf} supplies the Unicode
characters in the Lean listings.

The check validates the formal statements under Lean's foundations and
the installed compiler and dependencies. It does not by itself establish
historical novelty or replace expert review of the translation from an
informal claim to its definitions. The public map and count definitions are
printed verbatim below so that this translation can be inspected directly.

\appendix
\section{Lean source: complete local dependency chain}
The following files are reproduced verbatim from the verified project, in
dependency order, starting from the public theorem and the counterexample
dichotomy. They compile as modules within that project, with its pinned
imports; the listings are not a single standalone Lean file. This appendix
includes every transitive local module needed by those two endpoints.
The much larger imported upstream and Mathlib libraries remain external.
\subsection{\texorpdfstring{\texttt{CollatzTargetDensity/Dynamics.lean}}{CollatzTargetDensity/Dynamics.lean}}
\begin{Verbatim}[fontsize=\footnotesize,breaklines=true,breakanywhere=true,numbers=left,numbersep=5pt]
import Mathlib.Logic.Function.Iterate
import Lean.Elab.Tactic.Omega

/-!+# Nonperiodic seed selection

Elementary deterministic dynamics used by the proposed all-target extension.
No Collatz convergence hypothesis occurs in these statements.
-/

namespace CollatzTargetDensity

def Nonperiodic {α : Type*} (f : α → α) (x : α) : Prop :=
  ∀ k : ℕ, 0 < k → (f^[k]) x ≠ x

theorem periodic_multiples {α : Type*} {f : α → α} {x : α} {p : ℕ}
    (hp : (f^[p]) x = x) (q : ℕ) : (f^[q * p]) x = x := by
  induction q with
  | zero => simp
  | succ q ih => rw [Nat.succ_mul, Function.iterate_add_apply, hp, ih]

/-- The restriction of any function to its periodic points is injective. -/
theorem periodic_predecessors_eq {α : Type*} {f : α → α} {x y : α}
    {p q : ℕ} (hp : 0 < p) (hq : 0 < q)
    (hpx : (f^[p]) x = x) (hqy : (f^[q]) y = y)
    (hxy : f x = f y) : x = y := by
  have hprod : 0 < q * p := Nat.mul_pos hq hp
  have hx := periodic_multiples hpx q
  have hy : (f^[q * p]) y = y := by
    simpa only [Nat.mul_comm] using periodic_multiples hqy p
  have he : q * p = (q * p - 1) + 1 := by omega
  have hiter : (f^[q * p]) x = (f^[q * p]) y := by
    rw [he, Function.iterate_add_apply, Function.iterate_add_apply]
    simp only [Function.iterate_one, hxy]
  exact hx.symm.trans (hiter.trans hy)

/-- Two distinct immediate predecessors provide a nonperiodic choice. -/
theorem nonperiodic_choice {α : Type*} {f : α → α} {x y : α}
    (hne : x ≠ y) (hxy : f x = f y) : Nonperiodic f x ∨ Nonperiodic f y := by
  classical
  by_cases hx : Nonperiodic f x
  · exact Or.inl hx
  · right
    obtain ⟨p, hp, hpx⟩ : ∃ p, ∃ _ : 0 < p, (f^[p]) x = x := by
      simpa only [Nonperiodic, not_forall, not_not] using hx
    intro q hq hqy
    exact hne (periodic_predecessors_eq hp hq hpx hqy hxy)

/-- Repeated hits would make the target itself periodic. -/
theorem nonperiodic_hit_time_unique {α : Type*} {f : α → α} {source target : α}
    (ht : Nonperiodic f target) {d e : ℕ}
    (hd : (f^[d]) source = target) (he : (f^[e]) source = target) : d = e := by
  have hnot (a b : ℕ) (ha : (f^[a]) source = target)
      (hb : (f^[b]) source = target) (hab : a < b) : False := by
    have hsplit : b = (b - a) + a := by omega
    rw [hsplit, Function.iterate_add_apply, ha] at hb
    exact ht (b - a) (by omega) hb
  by_cases hde : d < e
  · exact False.elim (hnot d e hd he hde)
  · by_cases hed : e < d
    · exact False.elim (hnot e d he hd hed)
    · omega

/-- A residue-wise supply of two distinct children is enough; neither child's
    eventual convergence is required. The predicate records residue and size. -/
theorem select_nonperiodic_child {α : Type*} {f : α → α} (target : α)
    (P : α → Prop)
    (hchildren : ∃ x y, P x ∧ P y ∧ f x = target ∧ f y = target ∧ x ≠ y) :
    ∃ x, P x ∧ f x = target ∧ Nonperiodic f x := by
  obtain ⟨x, y, hx, hy, hfx, hfy, hne⟩ := hchildren
  rcases nonperiodic_choice hne (hfx.trans hfy.symm) with h | h
  · exact ⟨x, hx, hfx, h⟩
  · exact ⟨y, hy, hfy, h⟩

end CollatzTargetDensity
\end{Verbatim}
\subsection{\texorpdfstring{\texttt{CollatzTargetDensity/SeedSupply.lean}}{CollatzTargetDensity/SeedSupply.lean}}
\begin{Verbatim}[fontsize=\footnotesize,breaklines=true,breakanywhere=true,numbers=left,numbersep=5pt]
import CollatzTargetDensity.Dynamics
import Erdos1135.Tao.Syracuse.Basic

namespace CollatzTargetDensity
open Erdos1135

/-- The arithmetic input to the analytic argument: every finite ternary class
    contains arbitrarily large nonperiodic immediate predecessors. -/
def SeedSupply (target : ℕ) : Prop :=
  ∀ q X : ℕ, ∀ r : Fin (3 ^ q), ∃ M : ℕ,
    Odd M ∧ X < M ∧ Tao.syracuse M = target ∧ M % 3 ^ q = r ∧
      Nonperiodic Tao.syracuse M

end CollatzTargetDensity
\end{Verbatim}
\subsection{\texorpdfstring{\texttt{CollatzTargetDensity/ResidueSeedCore.lean}}{CollatzTargetDensity/ResidueSeedCore.lean}}
\begin{Verbatim}[fontsize=\footnotesize,breaklines=true,breakanywhere=true,numbers=left,numbersep=5pt]
import Erdos1135.ND.PositiveDensity.A5HistoricalAdjacentCylinderPreviousPhysicalLowerShellNoGo
import Mathlib.NumberTheory.Multiplicity

/-!
# Arithmetic residue coverage, isolated from density imports

The following proofs are copied without mathematical change from Lech Mazur's
Apache-2.0 source package, corrected commit
0aef8b0cfaaf6959500063472f2e328c60df8d54,
`SuccessfulRootTernaryResidueCoverage.lean`. Modification: namespace changed and
unneeded target-counting imports and later analytic declarations omitted.
This is known arithmetic, not a claimed new result.
-/
namespace CollatzTargetDensity.ResidueSeedCore
open Erdos1135 Erdos1135.ND.PositiveDensity

theorem ndCollapseSource_add (m n : ℕ) :
    ndA5PreviousPhysicalCollapseSource (m + n) =
      ndA5PreviousPhysicalCollapseSource m + 4 ^ m * ndA5PreviousPhysicalCollapseSource n := by
  have ha := three_mul_ndA5PreviousPhysicalCollapseSource_add_one (m + n)
  have hm := three_mul_ndA5PreviousPhysicalCollapseSource_add_one m
  have hn := three_mul_ndA5PreviousPhysicalCollapseSource_add_one n
  rw [pow_add, ← hm, ← hn] at ha
  nlinarith

theorem threePow_succ_dvd_fourPow_sub_one_iff (q n : ℕ) :
    3 ^ (q + 1) ∣ 4 ^ n - 1 ↔ 3 ^ q ∣ n := by
  rw [pow_dvd_iff_le_emultiplicity, pow_dvd_iff_le_emultiplicity]
  have h := Nat.emultiplicity_pow_sub_pow
    (p := 3) (x := 4) (y := 1)
    (by norm_num : Nat.Prime 3) (by norm_num : Odd 3)
    (by norm_num : 3 ∣ 4 - 1) (by norm_num : ¬ 3 ∣ 4) n
  simp only [one_pow] at h
  rw [h]
  norm_num [Nat.Prime.emultiplicity_self (by norm_num : Nat.Prime 3)]
  rw [add_comm (1 : ENat), ENat.add_le_add_iff_right (by simp)]

theorem threePow_dvd_collapseSource_iff (q n : ℕ) :
    3 ^ q ∣ ndA5PreviousPhysicalCollapseSource n ↔ 3 ^ q ∣ n := by
  have h := three_mul_ndA5PreviousPhysicalCollapseSource_add_one n
  have he : 4 ^ n - 1 = 3 * ndA5PreviousPhysicalCollapseSource n := by omega
  have hl := threePow_succ_dvd_fourPow_sub_one_iff q n
  rw [he, pow_succ, Nat.mul_comm (3 ^ q) 3,
    Nat.mul_dvd_mul_iff_left (by norm_num : 0 < 3)] at hl
  exact hl

theorem ndCollapseSource_modEq_of_add_iff (q m n : ℕ) :
    Nat.ModEq (3 ^ q) (ndA5PreviousPhysicalCollapseSource m)
      (ndA5PreviousPhysicalCollapseSource (m + n)) ↔ 3 ^ q ∣ n := by
  rw [ndCollapseSource_add, Nat.left_modEq_add_iff]
  have hc : Nat.Coprime (3 ^ q) (4 ^ m) := by
    exact (show Nat.Coprime 3 4 by norm_num).pow q m
  rw [hc.dvd_mul_left, threePow_dvd_collapseSource_iff]

theorem ndCollapseSource_residue_injective (q : ℕ) :
    Function.Injective (fun m : Fin (3 ^ q) =>
      (⟨ndA5PreviousPhysicalCollapseSource m % 3 ^ q,
        Nat.mod_lt _ (by positivity)⟩ : Fin (3 ^ q))) := by
  intro m n he
  apply Fin.ext
  have heq : Nat.ModEq (3 ^ q) (ndA5PreviousPhysicalCollapseSource m)
      (ndA5PreviousPhysicalCollapseSource n) := congrArg Fin.val he
  suffices h : ∀ (a b : Fin (3 ^ q)), a ≤ b →
      Nat.ModEq (3 ^ q) (ndA5PreviousPhysicalCollapseSource a)
        (ndA5PreviousPhysicalCollapseSource b) → (a : ℕ) = b by
    rcases le_total m n with hmn | hnm
    · exact h m n hmn heq
    · exact (h n m hnm heq.symm).symm
  intro a b hab hmod
  obtain ⟨d, hd⟩ := Nat.exists_eq_add_of_le (show (a : ℕ) ≤ b from hab)
  rw [hd] at hmod
  have hdiv := (ndCollapseSource_modEq_of_add_iff q a d).mp hmod
  have hdlt : d < 3 ^ q := by have h := b.isLt; omega
  have hz : d = 0 := Nat.eq_zero_of_dvd_of_lt hdiv hdlt
  omega

/-- The existing successful family is a permutation of every finite
ternary residue space. No finite modulus range is enumerated. -/
theorem ndCollapseSource_residue_surjective (q : ℕ) :
    Function.Surjective (fun m : Fin (3 ^ q) =>
      (⟨ndA5PreviousPhysicalCollapseSource m % 3 ^ q,
        Nat.mod_lt _ (by positivity)⟩ : Fin (3 ^ q))) :=
  Finite.surjective_of_injective (ndCollapseSource_residue_injective q)

/-- **ARBITRARILY HIGH SUCCESSFUL RESIDUE COVERAGE.** Every ternary class
contains unbounded odd integers that reach one in one Syracuse step. -/
theorem exists_large_oneStepSuccessful_root_in_residue
    (q X : ℕ) (r : Fin (3 ^ q)) :
    ∃ M : ℕ, X < M ∧ Odd M ∧ Tao.syracuse M = 1 ∧ M % 3 ^ q = r := by
  obtain ⟨m, hm⟩ := ndCollapseSource_residue_surjective q r
  let k := (m : ℕ) + (X + 1) * 3 ^ q
  have hQ : 1 ≤ 3 ^ q := by
    have h : 0 < 3 ^ q := by positivity
    omega
  have hXk : X < k := by
    have h := Nat.mul_le_mul_left (X + 1) hQ
    dsimp only [k]
    omega
  have hmod : Nat.ModEq (3 ^ q) (ndA5PreviousPhysicalCollapseSource m)
      (ndA5PreviousPhysicalCollapseSource k) :=
    (ndCollapseSource_modEq_of_add_iff q m ((X + 1) * 3 ^ q)).2
      (Nat.dvd_mul_left _ _)
  have hr : ndA5PreviousPhysicalCollapseSource k % 3 ^ q = r := by
    change ndA5PreviousPhysicalCollapseSource m % 3 ^ q =
      ndA5PreviousPhysicalCollapseSource k % 3 ^ q at hmod
    exact hmod.symm.trans (congrArg Fin.val hm)
  obtain ⟨j, hj⟩ := Nat.exists_eq_succ_of_ne_zero (by omega : k ≠ 0)
  refine ⟨ndA5PreviousPhysicalCollapseSource k,
    hXk.trans_le (self_le_ndA5PreviousPhysicalCollapseSource k), ?_, ?_, hr⟩
  · rw [hj]
    exact ndA5PreviousPhysicalCollapseSource_odd j
  · rw [hj]
    exact syracuse_ndA5PreviousPhysicalCollapseSource_succ j


end CollatzTargetDensity.ResidueSeedCore
\end{Verbatim}
\subsection{\texorpdfstring{\texttt{CollatzTargetDensity/ArithmeticSeeds.lean}}{CollatzTargetDensity/ArithmeticSeeds.lean}}
\begin{Verbatim}[fontsize=\footnotesize,breaklines=true,breakanywhere=true,numbers=left,numbersep=5pt]
import CollatzTargetDensity.SeedSupply
import CollatzTargetDensity.ResidueSeedCore

/-!+# Arbitrary-target nonperiodic seed supply

The residue coverage is the classical backward-mixing phenomenon. The proof
below uses the imported successful-root coverage as an arithmetic parameter
space, then transports it to another target by an affine map. It does not
assert that the new target converges to 1.
-/

namespace CollatzTargetDensity
open Erdos1135 Erdos1135.ND.PositiveDensity
noncomputable section

theorem exists_odd_child {target : ℕ} (hodd : Odd target) (hthree : target % 3 ≠ 0) :
    ∃ c : ℕ, Odd c ∧ Tao.syracuse c = target := by
  have hnotdvd : ¬2 ∣ target := by
    intro hdiv
    exact (Nat.not_even_iff_odd.mpr hodd) (even_iff_two_dvd.mpr hdiv)
  by_cases h : target % 3 = 1
  · let c := (4 * target - 1) / 3
    have he : 3 * c + 1 = 4 * target := by dsimp [c]; omega
    refine ⟨c, ⟨c / 2, by omega⟩, ?_⟩
    rw [Tao.syracuse_eq_ordCompl_two, he]
    exact Nat.ordCompl_pow_mul_of_not_dvd 2 Nat.prime_two hnotdvd
  · let c := (2 * target - 1) / 3
    have he : 3 * c + 1 = 2 * target := by dsimp [c]; omega
    refine ⟨c, ⟨c / 2, by omega⟩, ?_⟩
    rw [Tao.syracuse_eq_ordCompl_two, he]
    exact Nat.ordCompl_pow_mul_of_not_dvd 1 Nat.prime_two hnotdvd

/-- The affine transform x ↦ (3c+1)x+c takes one-step predecessors of 1
    to one-step predecessors of S(c), and permutes every ternary residue space. -/
theorem large_child_in_residue {target c : ℕ} (hc : Odd c)
    (hct : Tao.syracuse c = target) (q X : ℕ) (r : Fin (3 ^ q)) :
    ∃ M : ℕ, Odd M ∧ X < M ∧ Tao.syracuse M = target ∧ M % 3 ^ q = r := by
  have hcop : Nat.Coprime (3 * c + 1) 3 := by
    rw [Nat.add_comm, Nat.coprime_add_mul_left_left]
    norm_num
  have hunit : IsUnit ((3 * c + 1 : ℕ) : ZMod (3 ^ q)) := by
    apply (ZMod.isUnit_iff_coprime _ _).2
    simpa using hcop.pow 1 q
  let f : ZMod (3 ^ q) → ZMod (3 ^ q) :=
    fun x => ((3 * c + 1 : ℕ) : ZMod (3 ^ q)) * x + c
  have hinj : Function.Injective f := by
    intro x y hxy
    exact hunit.mul_left_cancel (add_right_cancel hxy)
  obtain ⟨y, hy⟩ := Finite.surjective_of_injective hinj (r.val : ZMod (3 ^ q))
  obtain ⟨x, hx, _hxodd, hxhit, hxres⟩ := ResidueSeedCore.exists_large_oneStepSuccessful_root_in_residue
    q X ⟨y.val, ZMod.val_lt y⟩
  let M := (3 * c + 1) * x + c
  have hxcast : (x : ZMod (3 ^ q)) = y := by
    calc
      _ = ((x % 3 ^ q : ℕ) : ZMod (3 ^ q)) := by simp
      _ = ((y.val : ℕ) : ZMod (3 ^ q)) := congrArg (fun a : ℕ => (a : ZMod (3 ^ q))) hxres
      _ = y := ZMod.natCast_zmod_val y
  have hMcast : (M : ZMod (3 ^ q)) = (r.val : ZMod (3 ^ q)) := by
    simpa only [M, f, Nat.cast_add, Nat.cast_mul, Nat.cast_ofNat, hxcast] using hy
  refine ⟨M, ?_, ?_, ?_, ?_⟩
  · obtain ⟨k, hk⟩ := hc
    refine ⟨(3 * k + 2) * x + k, ?_⟩
    dsimp [M]
    rw [hk]
    ring
  · dsimp [M]
    nlinarith
  · have he : 3 * M + 1 = (3 * c + 1) * (3 * x + 1) := by dsimp [M]; ring
    rw [Tao.syracuse_eq_ordCompl_two, he, Nat.ordCompl_mul,
      ← Tao.syracuse_eq_ordCompl_two, ← Tao.syracuse_eq_ordCompl_two,
      hct, hxhit, mul_one]
  · have hv := congrArg ZMod.val hMcast
    simpa only [ZMod.val_natCast, Nat.mod_eq_of_lt r.isLt] using hv

/-- No unproved dynamical assumption: two distinct children in the same
    residue class allow exclusion of the unique possible periodic child. -/
theorem seed_supply {target : ℕ} (hodd : Odd target) (hthree : target % 3 ≠ 0) :
    SeedSupply target := by
  obtain ⟨c, hc, hct⟩ := exists_odd_child hodd hthree
  intro q X r
  obtain ⟨x, hxodd, hxlarge, hxhit, hxres⟩ := large_child_in_residue hc hct q X r
  obtain ⟨y, hyodd, hylarge, hyhit, hyres⟩ := large_child_in_residue hc hct q x r
  rcases nonperiodic_choice (Nat.ne_of_lt hylarge) (hxhit.trans hyhit.symm) with hx | hy
  · exact ⟨x, hxodd, hxlarge, hxhit, hxres, hx⟩
  · exact ⟨y, hyodd, hxlarge.trans hylarge, hyhit, hyres, hy⟩

end
end CollatzTargetDensity
\end{Verbatim}
\subsection{\texorpdfstring{\texttt{CollatzTargetDensity/SourceCharge.lean}}{CollatzTargetDensity/SourceCharge.lean}}
\begin{Verbatim}[fontsize=\footnotesize,breaklines=true,breakanywhere=true,numbers=left,numbersep=5pt]
import CollatzTargetDensity.Dynamics
import Erdos1135.ND.PositiveDensity.Geom2ShiftedWideSymmetricTerminalWeightedCensus

/-!+# Source charging for nonperiodic seeds

Adapted from Lech Mazur's Apache-2.0 positive-density source package, corrected
commit 0aef8b0cfaaf6959500063472f2e328c60df8d54. The counting inequalities and
their proofs are retained; the dynamical hypothesis is changed from convergence
to 1 to nonperiodicity. These adaptations are separate from the upstream files.
-/

namespace Erdos1135.ND.PositiveDensity
open scoped BigOperators
open CollatzTargetDensity
noncomputable section
namespace NDGeom2ShiftedWideSymmetricRootSideUniformFloorState

theorem fullTerminal_eq_of_nonperiodic_ancestor_source_eq
    {U : NDGeom2ShiftedWideSymmetricRootSideUniformFloorState}
    {cap : ℕ → ℕ} {n : ℕ} {shift : (U.forwardIterate cap n).state.Label → ℕ} {K : ℕ}
    (hseed : ∀ i, Nonperiodic Tao.syracuse (U.state.root i))
    {z w : U.FullTerminalAt cap n shift K}
    (ha : fullTerminalAncestor z = fullTerminalAncestor w)
    (hs : ndGeom2PredictableRootSideBoundedOvershootIncidenceSource z =
      ndGeom2PredictableRootSideBoundedOvershootIncidenceSource w) : z = w := by
  let p := U.fullTerminalPath cap n shift K z
  let q := U.fullTerminalPath cap n shift K w
  have hq : (Tao.syracuse^[q.depth])
      (ndGeom2PredictableRootSideBoundedOvershootIncidenceSource z) =
      U.state.root (fullTerminalAncestor z) := by rw [hs, ha]; exact q.terminal_eq
  have hd := nonperiodic_hit_time_unique (hseed _) p.terminal_eq hq
  have hw := NDGeom2RootSideSyracusePath.word_eq_of_source_depth_eq p q hs hd
  have hpw := U.fullTerminalPath_word_eq_reverse cap n shift K z
  have hqw := U.fullTerminalPath_word_eq_reverse cap n shift K w
  apply fullTerminal_eq_of_ancestor_word_eq ha
  apply List.reverse_injective
  exact hpw.symm.trans (hw.trans hqw)

theorem fullTerminal_mass_mul_lower_le_frozenPotential_mul_card_of_nonperiodic
    (U : NDGeom2ShiftedWideSymmetricRootSideUniformFloorState)
    (cap : ℕ → ℕ) (n : ℕ)
    (shift : (U.forwardIterate cap n).state.Label → ℕ) (K : ℕ)
    (hseed : ∀ i, Nonperiodic Tao.syracuse (U.state.root i))
    (X : ℝ)
    (hX : ∀ z : U.FullTerminalAt cap n shift K,
      X ≤ (ndGeom2PredictableRootSideBoundedOvershootIncidenceSource z : ℝ)) :
    X * U.fullTerminalMass cap n shift K ≤
      U.parentSourcePotential * (U.fullTerminalSources cap n shift K).card := by
  classical
  letI := U.state.labelFintype
  letI := (U.forwardIterate cap n).state.labelFintype
  let W := fun z : U.FullTerminalAt cap n shift K =>
    ndGeom2PredictableRootSideBoundedOvershootIncidenceWeight
      (U.forwardIterate cap n).state.outerWeight z
  let P := fun i : U.state.Label => (U.state.root i : ℝ) * U.state.outerWeight i
  let cell := fun z : U.FullTerminalAt cap n shift K =>
    (fullTerminalAncestor z, ndGeom2PredictableRootSideBoundedOvershootIncidenceSource z)
  let sources := U.fullTerminalSources cap n shift K
  have hinj : Function.Injective cell := by
    intro z w h
    exact fullTerminal_eq_of_nonperiodic_ancestor_source_eq hseed
      (congrArg Prod.fst h) (congrArg Prod.snd h)
  have hsub : Finset.univ.image cell ⊆ Finset.univ.product sources := by
    intro c hc
    obtain ⟨z, hz, rfl⟩ := Finset.mem_image.mp hc
    exact Finset.mem_product.mpr ⟨Finset.mem_univ _, Finset.mem_image.mpr ⟨z, hz, rfl⟩⟩
  calc
    X * U.fullTerminalMass cap n shift K = ∑ z, X * W z := by
      change X * (∑ z, W z) = _
      rw [Finset.mul_sum]
    _ ≤ ∑ z : U.FullTerminalAt cap n shift K, P (fullTerminalAncestor z) := by
      apply Finset.sum_le_sum
      intro z _
      have hw := ndGeom2PredictableRootSideBoundedOvershootIncidenceWeight_nonneg
        (U.forwardIterate cap n).state.outerWeight
        (U.forwardIterate cap n).state.weight_nonneg z
      exact (mul_le_mul_of_nonneg_right (hX z) hw).trans
        (U.fullTerminal_source_mul_weight_le_frozenOwner cap n shift K z)
    _ = ∑ c ∈ Finset.univ.image cell, P c.1 := by
      rw [Finset.sum_image]
      exact fun a _ b _ h => hinj h
    _ ≤ ∑ c ∈ Finset.univ.product sources, P c.1 := by
      apply Finset.sum_le_sum_of_subset_of_nonneg hsub
      intro c _ _
      exact mul_nonneg (Nat.cast_nonneg _) (U.state.weight_nonneg c.1)
    _ = U.parentSourcePotential * (U.fullTerminalSources cap n shift K).card := by
      calc
        _ = ∑ i : U.state.Label, ∑ _s ∈ sources, P i := Finset.sum_product _ _ _
        _ = _ := ?_
      simp only [Finset.sum_const, nsmul_eq_mul, ← Finset.mul_sum]
      unfold parentSourcePotential ndGeom2PredictableRootSideParentSourcePotential
        NDGeom2ShiftedWideSymmetricRegenerativeState.rootSideUniformFloorAllLabels
      dsimp only [P, sources]
      ring


theorem fullTerminal_weighted_source_charge_of_nonperiodic
    (U : NDGeom2ShiftedWideSymmetricRootSideUniformFloorState)
    (cap : ℕ → ℕ) (n : ℕ)
    (shift : (U.forwardIterate cap n).state.Label → ℕ) (K : ℕ)
    (hseed : ∀ i, Nonperiodic Tao.syracuse (U.state.root i))
    (X : ℝ) (hX : 0 < X)
    (hsource : ∀ z : U.FullTerminalAt cap n shift K,
      X ≤ (ndGeom2PredictableRootSideBoundedOvershootIncidenceSource z : ℝ))
    (psi : ℕ → ℝ) (hp : ∀ x, 0 ≤ psi x) :
    letI := (U.forwardIterate cap n).state.labelFintype
    (∑ z : U.FullTerminalAt cap n shift K,
      ndGeom2PredictableRootSideBoundedOvershootIncidenceWeight
        (U.forwardIterate cap n).state.outerWeight z *
        psi (ndGeom2PredictableRootSideBoundedOvershootIncidenceSource z)) ≤
      (U.parentSourcePotential / X) *
        ∑ x ∈ U.fullTerminalSources cap n shift K, psi x := by
  classical
  letI := U.state.labelFintype
  letI := (U.forwardIterate cap n).state.labelFintype
  let W := fun z : U.FullTerminalAt cap n shift K =>
    ndGeom2PredictableRootSideBoundedOvershootIncidenceWeight
      (U.forwardIterate cap n).state.outerWeight z
  let P := fun i : U.state.Label => (U.state.root i : ℝ) * U.state.outerWeight i
  let cell := fun z : U.FullTerminalAt cap n shift K =>
    (fullTerminalAncestor z, ndGeom2PredictableRootSideBoundedOvershootIncidenceSource z)
  let sources := U.fullTerminalSources cap n shift K
  have hinj : Function.Injective cell := by
    intro z w h
    exact fullTerminal_eq_of_nonperiodic_ancestor_source_eq hseed
      (congrArg Prod.fst h) (congrArg Prod.snd h)
  have hsub : Finset.univ.image cell ⊆ Finset.univ.product sources := by
    intro c hc
    obtain ⟨z, hz, rfl⟩ := Finset.mem_image.mp hc
    exact Finset.mem_product.mpr ⟨Finset.mem_univ _, Finset.mem_image.mpr ⟨z, hz, rfl⟩⟩
  have hmul : X * (∑ z, W z * psi (cell z).2) ≤
      U.parentSourcePotential * ∑ x ∈ sources, psi x := by
    calc
      _ = ∑ z, (X * W z) * psi (cell z).2 := by rw [Finset.mul_sum]; simp [mul_assoc]
      _ ≤ ∑ z, P (cell z).1 * psi (cell z).2 := by
        apply Finset.sum_le_sum
        intro z _
        apply mul_le_mul_of_nonneg_right _ (hp _)
        exact (mul_le_mul_of_nonneg_right (hsource z)
          (ndGeom2PredictableRootSideBoundedOvershootIncidenceWeight_nonneg _
            (U.forwardIterate cap n).state.weight_nonneg z)).trans
          (U.fullTerminal_source_mul_weight_le_frozenOwner cap n shift K z)
      _ = ∑ c ∈ Finset.univ.image cell, P c.1 * psi c.2 := by
        rw [Finset.sum_image]; exact fun a _ b _ h => hinj h
      _ ≤ ∑ c ∈ Finset.univ.product sources, P c.1 * psi c.2 := by
        apply Finset.sum_le_sum_of_subset_of_nonneg hsub
        intro c _ _
        exact mul_nonneg (mul_nonneg (Nat.cast_nonneg _) (U.state.weight_nonneg c.1)) (hp _)
      _ = _ := by
        calc
          _ = ∑ i : U.state.Label, ∑ x ∈ sources, P i * psi x := Finset.sum_product _ _ _
          _ = _ := by
            simp only [← Finset.mul_sum, ← Finset.sum_mul]
            unfold parentSourcePotential ndGeom2PredictableRootSideParentSourcePotential
              NDGeom2ShiftedWideSymmetricRegenerativeState.rootSideUniformFloorAllLabels
            rfl
  have hdiv := (le_div_iff₀ hX).2 (by simpa [mul_comm] using hmul)
  simpa [W, cell, sources, div_eq_mul_inv, mul_comm, mul_left_comm, mul_assoc] using hdiv


theorem fullTerminal_weighted_residue_census_of_nonperiodic
    (U : NDGeom2ShiftedWideSymmetricRootSideUniformFloorState)
    (cap : ℕ → ℕ) (n : ℕ)
    (shift : (U.forwardIterate cap n).state.Label → ℕ) (K : ℕ)
    (hseed : ∀ i, Nonperiodic Tao.syracuse (U.state.root i))
    (q : ℕ) (X : ℝ) (hX : 0 < X) (hQ : ((3 ^ q : ℕ) : ℝ) ≤ X)
    (hsource : ∀ z : U.FullTerminalAt cap n shift K,
      X ≤ (ndGeom2PredictableRootSideBoundedOvershootIncidenceSource z : ℝ) ∧
      (ndGeom2PredictableRootSideBoundedOvershootIncidenceSource z : ℝ) < 32 * X)
    (psi : ZMod (3 ^ q) → ℝ) (hp : ∀ y, 0 ≤ psi y) :
    letI := (U.forwardIterate cap n).state.labelFintype
    (∑ z : U.FullTerminalAt cap n shift K,
      ndGeom2PredictableRootSideBoundedOvershootIncidenceWeight
        (U.forwardIterate cap n).state.outerWeight z *
        psi (ndGeom2PredictableRootSideBoundedOvershootIncidenceSource z : ZMod (3 ^ q))) ≤
      33 * U.parentSourcePotential * ndTernaryUniformMean q psi := by
  classical
  letI := (U.forwardIterate cap n).state.labelFintype
  have hc := U.fullTerminal_weighted_source_charge_of_nonperiodic cap n shift K hseed
    X hX (fun z => (hsource z).1) (fun x => psi (x : ZMod (3 ^ q))) (fun x => hp _)
  have hs := terminal_source_residue_sum_le (U.fullTerminalSources cap n shift K) q X hX hQ
    (fun x hx => by obtain ⟨z, _, rfl⟩ := Finset.mem_image.mp hx; exact (hsource z).2) psi hp
  calc
    _ ≤ _ := hc
    _ ≤ (U.parentSourcePotential / X) * (33 * X * ndTernaryUniformMean q psi) :=
      mul_le_mul_of_nonneg_left hs (div_nonneg U.frozenSourcePotential_nonneg hX.le)
    _ = _ := by field_simp


theorem fullTerminal_fan_two_conductor_bound_of_nonperiodic
    (U : NDGeom2ShiftedWideSymmetricRootSideUniformFloorState)
    (cap : ℕ → ℕ) (n : ℕ)
    (shift : (U.forwardIterate cap n).state.Label → ℕ) (K J : ℕ)
    (hseed : ∀ i, Nonperiodic Tao.syracuse (U.state.root i))
    (m k : ℕ) (hmk : m ≤ k) (X : ℝ) (hX : 0 < X)
    (hQ : ((3 ^ k : ℕ) : ℝ) ≤ X)
    (hsource : ∀ z : U.FullTerminalAt cap n shift K,
      X ≤ (ndGeom2PredictableRootSideBoundedOvershootIncidenceSource z : ℝ) ∧
      (ndGeom2PredictableRootSideBoundedOvershootIncidenceSource z : ℝ) < 32 * X)
    (epsilon : ℝ) (he : 0 ≤ epsilon)
    (hL1 : ndTernaryUniformMean k (fun y => |ndSyracuseUnitReferenceDensity k y -
      ndSyracuseUnitReferenceDensity m (Tao.taoZModThreeProjection hmk y)|) ≤ epsilon) :
    letI := (U.forwardIterate cap n).state.labelFintype
    (∑ z : U.FullTerminalAt cap n shift K,
      ndGeom2PredictableRootSideBoundedOvershootIncidenceWeight
        (U.forwardIterate cap n).state.outerWeight z *
        ∑ j : Fin J, (1 / 4 : ℝ) ^ j.val * ndSyracuseUnitReferenceDensity k
          (ndTerminalSourceFan j.val
            (ndGeom2PredictableRootSideBoundedOvershootIncidenceSource z) : ZMod (3 ^ k))) ≤
      (8 / 9 : ℝ) * ((3 ^ m : ℕ) : ℝ) * U.fullTerminalMass cap n shift K +
        44 * U.parentSourcePotential * epsilon := by
  classical
  letI := (U.forwardIterate cap n).state.labelFintype
  let W := fun z : U.FullTerminalAt cap n shift K =>
    ndGeom2PredictableRootSideBoundedOvershootIncidenceWeight
      (U.forwardIterate cap n).state.outerWeight z
  let source := fun z : U.FullTerminalAt cap n shift K =>
    ndGeom2PredictableRootSideBoundedOvershootIncidenceSource z
  let H := (2 / 3 : ℝ) * ((3 ^ m : ℕ) : ℝ)
  let delta := fun y : ZMod (3 ^ k) => |ndSyracuseUnitReferenceDensity k y -
    ndSyracuseUnitReferenceDensity m (Tao.taoZModThreeProjection hmk y)|
  have hW : ∀ z, 0 ≤ W z := fun z =>
    ndGeom2PredictableRootSideBoundedOvershootIncidenceWeight_nonneg _
      (U.forwardIterate cap n).state.weight_nonneg z
  have hmass : 0 ≤ U.fullTerminalMass cap n shift K := Finset.sum_nonneg fun z _ => hW z
  have hstep (j : ℕ) :
      (∑ z, W z * ndSyracuseUnitReferenceDensity k
        (ndTerminalSourceFan j (source z) : ZMod (3 ^ k))) ≤
      H * U.fullTerminalMass cap n shift K + 33 * U.parentSourcePotential * epsilon := by
    let psi := fun y : ZMod (3 ^ k) => delta (ndTerminalResidueFan k j y)
    have hp : ∀ y, 0 ≤ psi y := fun _ => abs_nonneg _
    have hmean : ndTernaryUniformMean k psi ≤ epsilon := by
      rw [show psi = (fun y => delta (ndTerminalResidueFan k j y)) from rfl,
        terminalResidueFan_fullMean]
      exact hL1
    have hc := U.fullTerminal_weighted_residue_census_of_nonperiodic cap n shift K hseed
      k X hX hQ hsource psi hp
    have hc' : (∑ z, W z * psi (source z : ZMod (3 ^ k))) ≤
        33 * U.parentSourcePotential * epsilon := hc.trans
      (mul_le_mul_of_nonneg_left hmean (mul_nonneg (by norm_num) U.frozenSourcePotential_nonneg))
    have hpoint (x : ℕ) : ndSyracuseUnitReferenceDensity k
        (ndTerminalSourceFan j x : ZMod (3 ^ k)) ≤ H + psi (x : ZMod (3 ^ k)) := by
      dsimp [psi, delta]
      rw [terminalResidueFan_natCast]
      have hb := unitReferenceDensity_finite_upper m
        (Tao.taoZModThreeProjection hmk (ndTerminalSourceFan j x : ZMod (3 ^ k)))
      have ha := le_abs_self (ndSyracuseUnitReferenceDensity k
          (ndTerminalSourceFan j x : ZMod (3 ^ k)) -
        ndSyracuseUnitReferenceDensity m
          (Tao.taoZModThreeProjection hmk (ndTerminalSourceFan j x : ZMod (3 ^ k))))
      dsimp [H]
      linarith
    calc
      _ ≤ ∑ z, W z * (H + psi (source z : ZMod (3 ^ k))) :=
        Finset.sum_le_sum fun z _ => mul_le_mul_of_nonneg_left (hpoint _) (hW z)
      _ = H * U.fullTerminalMass cap n shift K + ∑ z, W z * psi (source z : ZMod (3 ^ k)) := by
        simp_rw [mul_add]
        rw [Finset.sum_add_distrib, ← Finset.sum_mul]
        change (∑ z, W z) * H + _ = H * (∑ z, W z) + _
        ring
      _ ≤ _ := add_le_add_right hc' _
  let B := H * U.fullTerminalMass cap n shift K + 33 * U.parentSourcePotential * epsilon
  have hB : 0 ≤ B := add_nonneg (mul_nonneg (by dsimp [H]; positivity) hmass)
    (mul_nonneg (mul_nonneg (by norm_num) U.frozenSourcePotential_nonneg) he)
  calc
    _ = ∑ j : Fin J, (1 / 4 : ℝ) ^ j.val *
        ∑ z, W z * ndSyracuseUnitReferenceDensity k
          (ndTerminalSourceFan j.val (source z) : ZMod (3 ^ k)) := by
      simp_rw [Finset.mul_sum]
      rw [Finset.sum_comm]
      apply Finset.sum_congr rfl
      intro j _
      apply Finset.sum_congr rfl
      intro z _
      dsimp [W, source]
      ring
    _ ≤ ∑ j : Fin J, (1 / 4 : ℝ) ^ j.val * B :=
      Finset.sum_le_sum fun j _ => mul_le_mul_of_nonneg_left (hstep _) (by positivity)
    _ = (∑ j : Fin J, (1 / 4 : ℝ) ^ j.val) * B := (Finset.sum_mul _ _ _).symm
    _ ≤ (4 / 3 : ℝ) * B := mul_le_mul_of_nonneg_right (terminal_fan_coeff_sum_le J) hB
    _ = _ := by dsimp [B, H]; ring


theorem forwardCoreTerminalUnitMass_le_two_conductor_of_nonperiodic
    (U : NDGeom2ShiftedWideSymmetricRootSideUniformFloorState)
    (cap width : ℕ → ℕ) (n K m k : ℕ) (hmk : m ≤ k)
    (hseed : ∀ i, Nonperiodic Tao.syracuse (U.state.root i))
    (X : ℝ) (hX : 0 < X) (hQ : ((3 ^ k : ℕ) : ℝ) ≤ X)
    (hi : ∀ i : (U.forwardIterate cap n).state.Label,
      ndGeom2ShiftedWideSymmetricPhysicalIntervalMin (U.forwardIterate cap n).floor
        ((U.forwardIterate cap n).state.root i) < X ∧
      X ≤ ndGeom2ShiftedWideSymmetricPhysicalIntervalMax (U.forwardIterate cap n).floor
        ((U.forwardIterate cap n).state.root i))
    (hsource : ∀ z : (U.forwardIterate cap n).FullTerminalIncidence X hX hi,
      X ≤ (ndGeom2PredictableRootSideBoundedOvershootIncidenceSource z : ℝ) ∧
      (ndGeom2PredictableRootSideBoundedOvershootIncidenceSource z : ℝ) < 32 * X)
    (epsilon : ℝ) (he : 0 ≤ epsilon)
    (hL1 : ndTernaryUniformMean k (fun y => |ndSyracuseUnitReferenceDensity k y -
      ndSyracuseUnitReferenceDensity m (Tao.taoZModThreeProjection hmk y)|) ≤ epsilon) :
    U.forwardCoreTerminalUnitMass cap width n K k X hX hi ≤
      (8 / 9 : ℝ) * ((3 ^ m : ℕ) : ℝ) * U.fullTerminalMass cap n
        ((U.forwardIterate cap n).fullTerminalShift X hX hi) 1 +
        44 * U.parentSourcePotential * epsilon := by
  classical
  let V := U.forwardIterate cap n
  letI := V.state.labelFintype
  have hmask (i : V.state.Label) : U.forwardCoreOuterWeight cap width n i ≤ V.state.outerWeight i := by
    unfold forwardCoreOuterWeight
    split_ifs
    · exact le_rfl
    · exact V.state.weight_nonneg i
  apply (U.forwardCoreTerminalUnitMass_le_full_cap_one_fan cap width n K k X hX hi).trans
  apply le_trans _ (U.fullTerminal_fan_two_conductor_bound_of_nonperiodic cap n (V.fullTerminalShift X hX hi)
    1 (K / 2 + 1) hseed m k hmk X hX hQ hsource epsilon he hL1)
  apply Finset.sum_le_sum
  intro z _
  apply mul_le_mul_of_nonneg_right
  · unfold ndGeom2PredictableRootSideBoundedOvershootIncidenceWeight
    exact mul_le_mul_of_nonneg_right (hmask _)
      (ndGeom2PredictableRootSideBoundedOvershootIncidenceAtom_nonneg z)
  · exact Finset.sum_nonneg fun j _ => mul_nonneg (by positivity)
      (ndSyracuseUnitReferenceDensity_nonneg _ _)

/-- The physical shell estimate is independent of the eventual target. -/
theorem fullTerminal_source_window
    (U : NDGeom2ShiftedWideSymmetricRootSideUniformFloorState)
    (X : ℝ) (hX : 0 < X)
    (hi : ∀ i : U.state.Label,
      ndGeom2ShiftedWideSymmetricPhysicalIntervalMin U.floor (U.state.root i) < X ∧
      X ≤ ndGeom2ShiftedWideSymmetricPhysicalIntervalMax U.floor (U.state.root i))
    (z : U.FullTerminalIncidence X hX hi) :
    X ≤ (ndGeom2PredictableRootSideBoundedOvershootIncidenceSource z : ℝ) ∧
      (ndGeom2PredictableRootSideBoundedOvershootIncidenceSource z : ℝ) < 32 * X := by
  have hpacket : ∀ i, 16 ^ U.floor ≤ U.state.root i := fun i =>
    (Nat.pow_le_pow_right (by norm_num) (U.floor_le_base i)).trans (U.state.rootLower i)
  have hs := U.fullTerminalShift_spec X hX hi
    (ndGeom2PredictableRootSideBoundedOvershootIncidenceLabel z)
  have hb := ndGeom2RootSideCommonFloor_capOne_source_bounds U.state.root_odd
    (by have h := U.floor_twoHundred; omega) hpacket z
  exact ⟨hs.2.1.trans hb.1, by nlinarith only [hb.2, hs.2.2]⟩

theorem forwardCoreTerminalUnitMass_le_native_two_conductor_rate_of_nonperiodic
    (A : ℕ) (hA : 0 < A) :
    ∃ Cmix : ℝ, 0 ≤ Cmix ∧
      ∀ (U : NDGeom2ShiftedWideSymmetricRootSideUniformFloorState)
        (cap width : ℕ → ℕ) (n K m k : ℕ),
        1 ≤ m → (hmk : m ≤ k) → k ≤ (U.forwardIterate cap n).floor →
        (∀ i, Nonperiodic Tao.syracuse (U.state.root i)) →
        ∀ (X : ℝ) (hX : 0 < X)
          (hi : ∀ i : (U.forwardIterate cap n).state.Label,
            ndGeom2ShiftedWideSymmetricPhysicalIntervalMin (U.forwardIterate cap n).floor
              ((U.forwardIterate cap n).state.root i) < X ∧
            X ≤ ndGeom2ShiftedWideSymmetricPhysicalIntervalMax (U.forwardIterate cap n).floor
              ((U.forwardIterate cap n).state.root i)),
          0 < U.forwardCoreTerminalUnitMass cap width n K k X hX hi →
          U.forwardCoreTerminalUnitMass cap width n K k X hX hi ≤
            (8 / 9 : ℝ) * ((3 ^ m : ℕ) : ℝ) * U.fullTerminalMass cap n
              ((U.forwardIterate cap n).fullTerminalShift X hX hi) 1 +
              44 * U.parentSourcePotential * (Cmix / (m : ℝ) ^ A) := by
  obtain ⟨Cmix, hCmix, hmix⟩ := unitReferenceDensity_fullL1_polynomial_rate A hA
  refine ⟨Cmix, hCmix, ?_⟩
  intro U cap width n K m k hm hmk hk hseed X hX hi hT
  let V := U.forwardIterate cap n
  have hQ := U.terminal_conductor_room_of_pos_mark cap width n K k hk X hX hi hT
  exact U.forwardCoreTerminalUnitMass_le_two_conductor_of_nonperiodic cap width n K m k hmk hseed
    X hX hQ.le hi
    (fun z => V.fullTerminal_source_window X hX hi z)
    (Cmix / (m : ℝ) ^ A) (div_nonneg hCmix (by positivity)) (hmix m k hm hmk)


end NDGeom2ShiftedWideSymmetricRootSideUniformFloorState
end
end Erdos1135.ND.PositiveDensity
\end{Verbatim}
\subsection{\texorpdfstring{\texttt{CollatzTargetDensity/Activation.lean}}{CollatzTargetDensity/Activation.lean}}
\begin{Verbatim}[fontsize=\footnotesize,breaklines=true,breakanywhere=true,numbers=left,numbersep=5pt]
import CollatzTargetDensity.SeedSupply
import Erdos1135.ND.PositiveDensity.Geom2ShiftedWideSymmetricTerminalShiftImageRate

/-!+# Target-independent activation from residue-wise nonperiodic seeds

Adapted from Lech Mazur's Apache-2.0 positive-density package, corrected commit
0aef8b0cfaaf6959500063472f2e328c60df8d54. The original uniform analytic estimates
are imported unchanged. Only the physical seed-selection step is generalized.
-/

namespace CollatzTargetDensity
open Erdos1135 Erdos1135.ND.PositiveDensity
open Filter
open scoped Topology
noncomputable section

theorem seed_at_generation {target : ℕ} (hsupply : SeedSupply target)
    {b k : ℕ} (hb : 200 ≤ b) (hk : 1 ≤ k) (cap width : ℕ → ℕ)
    (hwidth : ∀ B, b ≤ B → 2 ≤ width B) (n X : ℕ) :
    ∃ (M : ℕ) (hM : Odd M) (hl : 16 ^ b ≤ M),
      X < M ∧ Tao.syracuse M = target ∧ Nonperiodic Tao.syracuse M ∧
      (2 / 3 : ℝ) * ndRootCoreProbabilityProduct b cap width n ≤
        (ndRootCoreSingletonState b M hb hM hl).forwardCoreMarkedMass cap width n k (fun _ => 1) := by
  obtain ⟨y, _, hy⟩ := exists_unit_rootCoreBackwardMark_ge_product
    (by omega : 9 ≤ b) hk cap width hwidth n
  let q := ndRootCoreBackwardConductor b k n
  let r : Fin (3 ^ q) := ⟨y.val, ZMod.val_lt y⟩
  obtain ⟨M, hM, hlarge, hhit, hres, hnonperiodic⟩ := hsupply q (max X (16 ^ b)) r
  have hl : 16 ^ b ≤ M := (le_max_right _ _).trans hlarge.le
  have heq : (M : ZMod (3 ^ q)) = y := by
    calc
      _ = ((M % 3 ^ q : ℕ) : ZMod (3 ^ q)) := by simp
      _ = ((y.val : ℕ) : ZMod (3 ^ q)) := congrArg (fun a : ℕ => (a : ZMod (3 ^ q))) hres
      _ = y := ZMod.natCast_zmod_val y
  refine ⟨M, hM, hl, (le_max_left _ _).trans_lt hlarge, hhit, hnonperiodic, ?_⟩
  rw [rootCoreSingletonState_markedMass_eq_backwardMark _ _ _ _ _ _ _ _ _ hk]
  simpa only [heq] using hy

theorem positive_core_limit {target : ℕ} (hsupply : SeedSupply target)
    {b : ℕ} (hb : 32 ^ 5 ≤ b) (K0 X : ℕ) :
    ∃ (M : ℕ) (hM : Odd M) (hl : 16 ^ b ≤ M),
      X < M ∧ Tao.syracuse M = target ∧ Nonperiodic Tao.syracuse M ∧ ∃ ell : ℝ, 0 < ell ∧
        Tendsto ((ndRootCoreSingletonState b M (by omega) hM hl).coreMarkedSequence
          (fun n => K0 + n / 100) ndRootCoreWidth) atTop (𝓝 ell) := by
  let cap (n : ℕ) := K0 + n / 100
  have hb200 : 200 ≤ b := by omega
  have hwidth (B : ℕ) (hB : b ≤ B) : 2 ≤ ndRootCoreWidth B := by
    have h := (rootCoreWidth_guard (hb.trans hB)).1
    omega
  have hupper (n : ℕ) : cap n ≤ (K0 + 1) * (n + 1) := by
    have h := Nat.div_le_self n 100
    dsimp [cap]
    nlinarith
  have hlower (n : ℕ) : K0 + n / 100 ≤ cap n := le_rfl
  obtain ⟨M0, hM0, hl0, _, _, _, _⟩ :=
    seed_at_generation hsupply hb200 (by norm_num : 1 ≤ 1)
      cap ndRootCoreWidth hwidth 0 0
  let U0 := ndRootCoreSingletonState b M0 hb200 hM0 hl0
  obtain ⟨p, hp, hprod⟩ := U0.exists_pos_le_coreProbabilityProduct hb cap K0 hlower
  change ∀ n, p ≤ ndRootCoreProbabilityProduct b cap ndRootCoreWidth n at hprod
  obtain ⟨C, hC, hlimall⟩ :=
    NDGeom2ShiftedWideSymmetricRootSideUniformFloorState.exists_coreMarkedLimit_with_root_uniform_tail
  have ht : ∀ᶠ N in atTop, ndRootCoreVariationTail b (K0 + 1) K0 C N < p / 3 :=
    (tendsto_order.mp (tendsto_rootCoreVariationTail_zero b (K0 + 1) K0 C)).2 (p / 3) (by positivity)
  obtain ⟨N, hN⟩ := Filter.eventually_atTop.mp ht
  let k := (U0.forwardIterate cap N).floor / 4
  have hk : 1 ≤ k := by have h := (U0.forwardIterate cap N).floor_twoHundred; dsimp [k]; omega
  obtain ⟨M, hM, hl, hX, hhit, hnonperiodic, hmark⟩ :=
    seed_at_generation hsupply hb200 hk cap ndRootCoreWidth hwidth N X
  let U := ndRootCoreSingletonState b M hb200 hM hl
  have hfloor := U.forward_floor_eq_of_floor_eq U0 cap cap rfl N
  have hmarked : (2 / 3 : ℝ) * ndRootCoreProbabilityProduct b cap ndRootCoreWidth N ≤
      U.coreMarkedSequence cap ndRootCoreWidth N := by
    unfold NDGeom2ShiftedWideSymmetricRootSideUniformFloorState.coreMarkedSequence
    rw [hfloor]
    exact hmark
  obtain ⟨ell, _, hlim, htail⟩ := hlimall U hb cap (K0 + 1) K0 hupper hlower
  have hdev := (abs_le.mp (htail N)).1
  have hD : U.denominator = 1 := rootCoreSingletonState_denominator_eq_one b M hb200 hM hl
  rw [hD, one_mul] at hdev
  change -(ndRootCoreVariationTail b (K0 + 1) K0 C N) ≤
    ell - U.coreMarkedSequence cap ndRootCoreWidth N at hdev
  refine ⟨M, hM, hl, hX, hhit, hnonperiodic, ell, ?_, hlim⟩
  have hn := hN N le_rfl
  have hpn := hprod N
  linarith


theorem uniform_terminal_mark {target : ℕ} (hsupply : SeedSupply target)
    {b : ℕ} (hb : 32 ^ 5 ≤ b) (K0 R : ℕ) :
    ∃ (M : ℕ) (hM : Odd M) (hl : 16 ^ b ≤ M),
      R < M ∧ Tao.syracuse M = target ∧ Nonperiodic Tao.syracuse M ∧ ∃ a : ℝ, 0 < a ∧
      ∀ᶠ n in atTop,
        let U := ndRootCoreSingletonState b M (by omega) hM hl;
        let cap := fun n => K0 + n / 100;
        let V := U.forwardIterate cap n;
        ∀ (X : ℝ) (hX : 0 < X)
        (hi : ∀ i : V.state.Label,
          ndGeom2ShiftedWideSymmetricPhysicalIntervalMin V.floor (V.state.root i) < X ∧
          X ≤ ndGeom2ShiftedWideSymmetricPhysicalIntervalMax V.floor (V.state.root i)),
        a ≤ U.forwardCoreTerminalUnitMass cap ndRootCoreWidth n (cap n) (V.floor / 4) X hX hi := by
  obtain ⟨M, hM, hl, hR, hhit, hnonperiodic, ell, hell, hlim⟩ :=
    positive_core_limit hsupply hb K0 R
  let U := ndRootCoreSingletonState b M (by omega) hM hl
  have hspan : U.rootSpan 0 := by intro i j; change M ≤ 2 ^ 0 * M; simp
  have hupper (n : ℕ) : K0 + n / 100 ≤ (K0 + 1) * (n + 1) := by
    have h := Nat.div_le_self n 100
    nlinarith
  exact ⟨M, hM, hl, hR, hhit, hnonperiodic, ell / 2, by positivity,
    U.eventually_terminalUnitMass_ge_of_pos_core_limit _ 0 (K0 + 1) K0 hspan hupper (fun _ => le_rfl) hell hlim⟩

end
end CollatzTargetDensity
\end{Verbatim}
\subsection{\texorpdfstring{\texttt{CollatzTargetDensity/RawBridge.lean}}{CollatzTargetDensity/RawBridge.lean}}
\begin{Verbatim}[fontsize=\footnotesize,breaklines=true,breakanywhere=true,numbers=left,numbersep=5pt]
import CollatzTargetDensity.SeedSupply
import Erdos1135.Tao.Syracuse.CollatzBridge
import Erdos1135.Parity

/-! Deterministic reachability and transfer to even targets. -/
namespace CollatzTargetDensity
open Erdos1135
noncomputable section

def oddReaches (target : ℕ) : Set ℕ :=
  {n | Odd n ∧ ∃ k : ℕ, (Tao.syracuse^[k]) n = target}

/-- Choose an odd predecessor prime to 3 of any target prime to 3.
    This construction also works when the target is even. -/
theorem exists_fertile_raw_child {target : ℕ} (hthree : target % 3 ≠ 0) :
    ∃ c : ℕ, Odd c ∧ c % 3 ≠ 0 ∧ Reaches c target := by
  have hbase : ∃ c v : ℕ, Odd c ∧ 3 * c + 1 = 2 ^ v * target := by
    by_cases h : target % 3 = 1
    · let c := (4 * target - 1) / 3
      have he : 3 * c + 1 = 4 * target := by dsimp [c]; omega
      exact ⟨c, 2, ⟨c / 2, by omega⟩, he⟩
    · let c := (2 * target - 1) / 3
      have he : 3 * c + 1 = 2 * target := by dsimp [c]; omega
      exact ⟨c, 1, ⟨c / 2, by omega⟩, he⟩
  have hfertile : ∃ c v : ℕ, Odd c ∧ c % 3 ≠ 0 ∧ 3 * c + 1 = 2 ^ v * target := by
    obtain ⟨c, v, hc, he⟩ := hbase
    by_cases h : c % 3 ≠ 0
    · exact ⟨c, v, hc, h, he⟩
    · refine ⟨4 * c + 1, v + 2, ⟨2 * c, by omega⟩, by omega, ?_⟩
      calc
        _ = 4 * (3 * c + 1) := by ring
        _ = 4 * (2 ^ v * target) := by rw [he]
        _ = _ := by rw [pow_add]; ring
  obtain ⟨c, v, hc, hunit, he⟩ := hfertile
  refine ⟨c, hc, hunit, v + 1, ?_⟩
  rw [Function.iterate_add_apply, Function.iterate_one,
    collatzStep_eq_three_mul_add_one_of_odd hc, he]
  exact Tao.collatz_iterate_powTwo_mul v target

def rawReaches (target : ℕ) : Set ℕ :=
  {n | 0 < n ∧ Reaches n target}

theorem oddReaches_subset_rawReaches {seed target : ℕ} (hseed : Reaches seed target) :
    oddReaches seed ⊆ rawReaches target := by
  intro n hn
  obtain ⟨hnodd, k, hk⟩ := hn
  have hreach : Reaches n seed := by
    simpa only [hk] using Tao.reaches_syracuse_iterate hnodd k
  exact ⟨hnodd.pos, hreach.trans hseed⟩

end
end CollatzTargetDensity
\end{Verbatim}
\subsection{\texorpdfstring{\texttt{CollatzTargetDensity/Assembly.lean}}{CollatzTargetDensity/Assembly.lean}}
\begin{Verbatim}[fontsize=\footnotesize,breaklines=true,breakanywhere=true,numbers=left,numbersep=5pt]
import CollatzTargetDensity.SourceCharge
import CollatzTargetDensity.Activation
import CollatzTargetDensity.RawBridge
import Erdos1135.ND.PositiveDensity.Geom2ShiftedWideSymmetricPositiveDensityAssembly

/-!
# Density assembly with an arbitrary target

Adapted from Lech Mazur's Apache-2.0 source, corrected commit
0aef8b0cfaaf6959500063472f2e328c60df8d54. The analytic inputs and interval
synchronization are unchanged. The conclusion records a supplied target;
nonperiodicity replaces seed convergence in the counting step.
-/
namespace Erdos1135.ND.PositiveDensity
open scoped BigOperators
open Filter CollatzTargetDensity
noncomputable section
namespace NDGeom2ShiftedWideSymmetricRootSideUniformFloorState

theorem eventually_pos_fullTerminalMass_of_nonperiodic_mark
    (U : NDGeom2ShiftedWideSymmetricRootSideUniformFloorState)
    (cap width : ℕ → ℕ) {a : ℝ}
    (hseed : ∀ i, Nonperiodic Tao.syracuse (U.state.root i))
    (ha : 0 < a)
    (hmark : ∀ᶠ n in atTop,
      let V := U.forwardIterate cap n;
      ∀ (X : ℝ) (hX : 0 < X)
        (hi : ∀ i : V.state.Label,
          ndGeom2ShiftedWideSymmetricPhysicalIntervalMin V.floor (V.state.root i) < X ∧
          X ≤ ndGeom2ShiftedWideSymmetricPhysicalIntervalMax V.floor (V.state.root i)),
        a ≤ U.forwardCoreTerminalUnitMass cap width n (cap n) (V.floor / 4) X hX hi) :
    ∃ eta : ℝ, 0 < eta ∧ ∀ᶠ n in atTop,
      let V := U.forwardIterate cap n;
      ∀ (X : ℝ) (hX : 0 < X)
        (hi : ∀ i : V.state.Label,
          ndGeom2ShiftedWideSymmetricPhysicalIntervalMin V.floor (V.state.root i) < X ∧
          X ≤ ndGeom2ShiftedWideSymmetricPhysicalIntervalMax V.floor (V.state.root i)),
        eta ≤ U.fullTerminalMass cap n (V.fullTerminalShift X hX hi) 1 := by
  obtain ⟨Cmix, hCmix, hrate⟩ := forwardCoreTerminalUnitMass_le_native_two_conductor_rate_of_nonperiodic 2 (by decide)
  obtain ⟨m, hm⟩ := exists_nat_gt (max 1 (88 * U.parentSourcePotential * Cmix / a))
  have hm1 : (1 : ℝ) < m := (le_max_left _ _).trans_lt hm
  have hmNat : 1 ≤ m := by exact_mod_cast hm1.le
  have hmPos : (0 : ℝ) < m := by linarith
  have hcoarse : 44 * U.parentSourcePotential * (Cmix / (m : ℝ) ^ 2) ≤ a / 2 := by
    have hmBig : 88 * U.parentSourcePotential * Cmix / a < (m : ℝ) :=
      (le_max_right _ _).trans_lt hm
    have hmProd := (div_lt_iff₀ ha).1 hmBig
    have hmSq : (m : ℝ) ≤ (m : ℝ) ^ 2 := by nlinarith
    have hpay := mul_le_mul_of_nonneg_left hmSq ha.le
    rw [← mul_div_assoc]
    apply (div_le_iff₀ (by positivity : (0 : ℝ) < (m : ℝ) ^ 2)).2
    nlinarith
  let H := (8 / 9 : ℝ) * ((3 ^ m : ℕ) : ℝ)
  have hH : 0 < H := by dsimp [H]; positivity
  refine ⟨a / (2 * H), by positivity, ?_⟩
  filter_upwards [hmark, eventually_ge_atTop (2 * m)] with n hn hnlarge
  intro V X hX hi
  have hf : U.floor + 2 * n ≤ V.floor := by
    have h := U.floor_add_two_mul_le_geometricIntervalStoppedIterate_floor cap 0 n
    rw [U.geometricIntervalStoppedIterate_floor_eq_iterate cap 0 n] at h
    simpa only [V, U.forwardIterate_eq_iterate cap n] using h
  have hmk : m ≤ V.floor / 4 := by omega
  have hT := hn X hX hi
  have h := hrate U cap width n (cap n) m (V.floor / 4) hmNat hmk
    (Nat.div_le_self _ _) hseed X hX hi (ha.trans_le hT)
  apply (div_le_iff₀ (mul_pos (by norm_num) hH)).2
  change a ≤ U.fullTerminalMass cap n (V.fullTerminalShift X hX hi) 1 * (2 * H)
  dsimp [H]
  nlinarith [hcoarse]


theorem targetDensity_of_eventually_pos_fullTerminalMass
    (U : NDGeom2ShiftedWideSymmetricRootSideUniformFloorState)
    (cap : ℕ → ℕ) {e L : ℕ} (he : U.rootSpan e)
    (hcap : ∀ n, cap n ≤ L * (n + 1))
    {eta : ℝ} (target : ℕ)
    (hseed : ∀ i, Nonperiodic Tao.syracuse (U.state.root i))
    (htarget : ∀ i, ∃ k : ℕ, (Tao.syracuse^[k]) (U.state.root i) = target)
    (hP : 0 < U.parentSourcePotential) (heta : 0 < eta)
    (hmass : ∀ᶠ n in atTop,
      let V := U.forwardIterate cap n;
      ∀ (X : ℝ) (hX : 0 < X)
        (hi : ∀ i : V.state.Label,
          ndGeom2ShiftedWideSymmetricPhysicalIntervalMin V.floor (V.state.root i) < X ∧
          X ≤ ndGeom2ShiftedWideSymmetricPhysicalIntervalMax V.floor (V.state.root i)),
        eta ≤ U.fullTerminalMass cap n (V.fullTerminalShift X hX hi) 1) :
    ∀ᶠ Y : ℕ in atTop,
      (eta / (32 * U.parentSourcePotential)) * (Y : ℝ) ≤
        (Terras.natCount (oddReaches target) Y : ℝ) := by
  classical
  obtain ⟨N0, hN0⟩ := eventually_atTop.mp hmass
  let Nstar := max N0 (1000000000 * (e + L + 3))
  let V0 := U.iterate cap Nstar
  letI := V0.state.labelFintype
  let X0 := ∑ i : V0.state.Label,
    max 0 (ndGeom2ShiftedWideSymmetricPhysicalIntervalMax V0.floor (V0.state.root i))
  have hstart (X : ℝ) (hXX : X0 < X) : ∀ i : V0.state.Label,
      ndGeom2ShiftedWideSymmetricPhysicalIntervalMax V0.floor (V0.state.root i) < X := by
    intro i
    have hsingle := Finset.single_le_sum
      (f := fun j : V0.state.Label => max 0
        (ndGeom2ShiftedWideSymmetricPhysicalIntervalMax V0.floor (V0.state.root j)))
      (fun j _ => le_max_left _ _) (Finset.mem_univ i)
    exact ((le_max_right _ _).trans hsingle).trans_lt hXX
  apply eventually_atTop.mpr
  refine ⟨Nat.ceil (32 * (max X0 0 + 1)), ?_⟩
  intro Y hY
  have hYlarge : 32 * (max X0 0 + 1) ≤ (Y : ℝ) :=
    (Nat.le_ceil _).trans (by exact_mod_cast hY)
  let X := (Y : ℝ) / 32
  have hX : 0 < X := by dsimp [X]; have := le_max_right X0 0; linarith
  have hXX : X0 < X := by dsimp [X]; have := le_max_left X0 0; linarith
  obtain ⟨n, hn, hi⟩ := U.exists_unstopped_generation_all_intervals_contain he cap hcap
    Nstar (le_max_right _ _) X (hstart X hXX)
  have hif : ∀ i : (U.forwardIterate cap n).state.Label,
      ndGeom2ShiftedWideSymmetricPhysicalIntervalMin (U.forwardIterate cap n).floor
        ((U.forwardIterate cap n).state.root i) < X ∧
      X ≤ ndGeom2ShiftedWideSymmetricPhysicalIntervalMax (U.forwardIterate cap n).floor
        ((U.forwardIterate cap n).state.root i) := by
    rw [U.forwardIterate_eq_iterate cap n]
    exact hi
  let shift := (U.forwardIterate cap n).fullTerminalShift X hX hif
  let sources := U.fullTerminalSources cap n shift 1
  have hm : eta ≤ U.fullTerminalMass cap n shift 1 :=
    hN0 n ((le_max_left _ _).trans hn) X hX hif
  have hwindow (z : U.FullTerminalAt cap n shift 1) :=
    (U.forwardIterate cap n).fullTerminal_source_window X hX hif z
  have hcountMass := U.fullTerminal_mass_mul_lower_le_frozenPotential_mul_card_of_nonperiodic
    cap n shift 1 hseed X (fun z => (hwindow z).1)
  have hmember (source : ℕ) (hs : source ∈ sources) :
      source ∈ oddReaches target ∧ X ≤ (source : ℝ) ∧ (source : ℝ) < 32 * X := by
    obtain ⟨z, _, rfl⟩ := Finset.mem_image.mp hs
    let p := U.fullTerminalPath cap n shift 1 z
    obtain ⟨k, hk⟩ := htarget (fullTerminalAncestor z)
    refine ⟨⟨p.sourceOdd, k + p.depth, ?_⟩, hwindow z⟩
    rw [Function.iterate_add_apply, p.terminal_eq, hk]
  have hrange : sources ⊆ Finset.range Y := by
    intro source hs
    have h := (hmember source hs).2.2
    have hlt : (source : ℝ) < Y := by dsimp [X] at h; linarith
    exact Finset.mem_range.mpr (by exact_mod_cast hlt)
  have htargetS : (↑sources : Set ℕ) ⊆ oddReaches target :=
    fun source hs => (hmember source hs).1
  have hcard : sources.card ≤ Terras.natCount (oddReaches target) Y :=
    finset_card_le_natCount_of_subset_range hrange htargetS
  have hcount : X * U.fullTerminalMass cap n shift 1 ≤
      U.parentSourcePotential * (Terras.natCount (oddReaches target) Y : ℝ) :=
    hcountMass.trans (mul_le_mul_of_nonneg_left (by exact_mod_cast hcard) hP.le)
  have htotal := (mul_le_mul_of_nonneg_left hm hX.le).trans hcount
  rw [div_mul_eq_mul_div]
  apply (div_le_iff₀ (by positivity : (0 : ℝ) < 32 * U.parentSourcePotential)).2
  dsimp [X] at htotal
  nlinarith only [htotal]


end NDGeom2ShiftedWideSymmetricRootSideUniformFloorState
end
end Erdos1135.ND.PositiveDensity

namespace CollatzTargetDensity
open Erdos1135 Erdos1135.ND.PositiveDensity Filter
noncomputable section

/-- Analytic conclusion from residue-wise nonperiodic seed supply.
    The arithmetic supply is an explicit input, not an assumed Collatz result. -/
theorem positive_density_of_seed_supply {target : ℕ} (hsupply : SeedSupply target) :
    ∃ c : ℝ, 0 < c ∧ ∀ᶠ X : ℕ in atTop,
      c * (X : ℝ) ≤ (Terras.natCount (oddReaches target) X : ℝ) := by
  classical
  obtain ⟨M, hM, hl, hMpos, hhit, hnonperiodic, a, ha, hmark⟩ :=
    uniform_terminal_mark hsupply (b := 32 ^ 5) (by rfl) 0 0
  let U := ndRootCoreSingletonState (32 ^ 5) M (by norm_num) hM hl
  let cap := fun n : ℕ => 0 + n / 100
  have hseed : ∀ i, Nonperiodic Tao.syracuse (U.state.root i) := fun _ => hnonperiodic
  have htarget : ∀ i, ∃ k : ℕ, (Tao.syracuse^[k]) (U.state.root i) = target := by
    intro i
    exact ⟨1, hhit⟩
  have hspan : U.rootSpan 0 := by intro i j; change M ≤ 2 ^ 0 * M; simp
  have hcap : ∀ n, cap n ≤ 1 * (n + 1) := by intro n; dsimp [cap]; omega
  have hP : 0 < U.parentSourcePotential := by
    change 0 < ndGeom2PredictableRootSideParentSourcePotential (Finset.univ : Finset Unit)
      (fun _ => (1 : ℝ)) (fun _ => M)
    simpa [ndGeom2PredictableRootSideParentSourcePotential] using
      (show (0 : ℝ) < M by exact_mod_cast hMpos)
  obtain ⟨eta, heta, hmass⟩ := U.eventually_pos_fullTerminalMass_of_nonperiodic_mark
    cap ndRootCoreWidth hseed ha hmark
  exact ⟨eta / (32 * U.parentSourcePotential), by positivity,
    U.targetDensity_of_eventually_pos_fullTerminalMass cap hspan hcap
      target hseed htarget hP heta hmass⟩

end
end CollatzTargetDensity
\end{Verbatim}
\subsection{\texorpdfstring{\texttt{CollatzTargetDensity/AllTargets.lean}}{CollatzTargetDensity/AllTargets.lean}}
\begin{Verbatim}[fontsize=\footnotesize,breaklines=true,breakanywhere=true,numbers=left,numbersep=5pt]
import CollatzTargetDensity.ArithmeticSeeds
import CollatzTargetDensity.Assembly
import CollatzTargetDensity.RawBridge

/-!+# Proposed all-target predecessor density endpoint

This file must pass the complete build and an axiom audit before its theorem
is reported as verified. The new argument depends on the imported analytic
machinery of Lech Mazur's positive-density theorem.
-/

namespace CollatzTargetDensity
open Erdos1135 Filter
noncomputable section

theorem positive_density_odd_target {target : ℕ}
    (hodd : Odd target) (hthree : target % 3 ≠ 0) :
    ∃ c : ℝ, 0 < c ∧ ∀ᶠ X : ℕ in atTop,
      c * (X : ℝ) ≤ (Terras.natCount (oddReaches target) X : ℝ) :=
  positive_density_of_seed_supply (seed_supply hodd hthree)

/-- Every positive target prime to 3 has a predecessor set of positive lower
    natural density. This is not a statement that every integer reaches 1. -/
theorem positive_density_all_targets {target : ℕ} (hthree : target % 3 ≠ 0) :
    ∃ c : ℝ, 0 < c ∧ ∃ X₀ : ℕ, ∀ X : ℕ, X₀ ≤ X →
      c * (X : ℝ) ≤ (Terras.natCount (rawReaches target) X : ℝ) := by
  obtain ⟨seed, hodd, hunit, hhit⟩ := exists_fertile_raw_child hthree
  obtain ⟨c, hc, hcount⟩ := positive_density_odd_target hodd hunit
  obtain ⟨X₀, hX₀⟩ := eventually_atTop.mp hcount
  refine ⟨c, hc, X₀, ?_⟩
  intro X hX
  have hmono := Terras.natCount_mono (oddReaches_subset_rawReaches hhit) X
  exact (hX₀ X hX).trans (by exact_mod_cast hmono)

end
end CollatzTargetDensity
\end{Verbatim}
\subsection{\texorpdfstring{\texttt{CollatzTargetDensity/ThreeDivisible.lean}}{CollatzTargetDensity/ThreeDivisible.lean}}
\begin{Verbatim}[fontsize=\footnotesize,breaklines=true,breakanywhere=true,numbers=left,numbersep=5pt]
import CollatzTargetDensity.RawBridge
import Erdos1135.Terras.Density.NaturalDensity

/-! Exact predecessor structure for the excluded targets divisible by 3. -/
namespace CollatzTargetDensity
open Erdos1135

theorem reaches_three_divisible_iff {target n : ℕ} (hthree : target % 3 = 0) :
    Reaches n target ↔ ∃ k : ℕ, n = 2 ^ k * target := by
  constructor
  · rintro ⟨m, hm⟩
    induction m generalizing n with
    | zero => exact ⟨0, by simpa using hm⟩
    | succ m ih =>
      rw [Function.iterate_succ_apply] at hm
      obtain ⟨k, hk⟩ := ih hm
      have hstep : collatzStep n % 3 = 0 := by
        rw [hk, Nat.mul_mod, hthree]
        simp
      have he : Even n := by
        by_contra hne
        rw [collatzStep_eq_three_mul_add_one_of_not_even hne] at hstep
        omega
      have hmod := Nat.even_iff.mp he
      refine ⟨k + 1, ?_⟩
      calc
        n = 2 * (n / 2) := by omega
        _ = 2 * collatzStep n := by rw [collatzStep_eq_div_two_of_even he]
        _ = 2 * (2 ^ k * target) := by rw [hk]
        _ = 2 ^ (k + 1) * target := by rw [pow_succ]; ring
  · rintro ⟨k, rfl⟩
    exact Tao.reaches_powTwo_mul k target

theorem three_divisible_count_at_dyadic_cutoff {target : ℕ}
    (hpos : 0 < target) (hthree : target % 3 = 0) (k : ℕ) :
    Terras.natCount (rawReaches target) (2 ^ k * target) = k := by
  classical
  have hmono : StrictMono (fun j : ℕ => 2 ^ j * target) := by
    intro i j hij
    exact Nat.mul_lt_mul_of_pos_right
      ((Nat.pow_lt_pow_iff_right (by decide : 1 < (2 : ℕ))).2 hij) hpos
  have hfilter : (Finset.range (2 ^ k * target)).filter (fun n => n ∈ rawReaches target) =
      (Finset.range k).image (fun j => 2 ^ j * target) := by
    ext n
    simp only [Finset.mem_filter, Finset.mem_range, Finset.mem_image]
    constructor
    · rintro ⟨hnlt, hn⟩
      obtain ⟨j, hj⟩ := (reaches_three_divisible_iff hthree).1 hn.2
      refine ⟨j, ?_, hj.symm⟩
      rw [hj] at hnlt
      exact hmono.lt_iff_lt.mp hnlt
    · rintro ⟨j, hj, rfl⟩
      exact ⟨hmono hj, Nat.mul_pos (by positivity) hpos, Tao.reaches_powTwo_mul j target⟩
  unfold Terras.natCount
  rw [Nat.count_eq_card_filter_range, hfilter,
    Finset.card_image_of_injective _ hmono.injective, Finset.card_range]

theorem square_le_two_pow {k : ℕ} (hk : 4 ≤ k) : k ^ 2 ≤ 2 ^ k := by
  induction k, hk using Nat.le_induction with
  | base => decide
  | succ k hk ih =>
    rw [pow_succ (2 : ℕ) k]
    nlinarith

theorem no_positive_lower_density_three_divisible {target : ℕ}
    (hpos : 0 < target) (hthree : target % 3 = 0) :
    ¬∃ c : ℝ, 0 < c ∧ ∃ X₀ : ℕ, ∀ X : ℕ, X₀ ≤ X →
      c * (X : ℝ) ≤ (Terras.natCount (rawReaches target) X : ℝ) := by
  rintro ⟨c, hc, X₀, hX₀⟩
  obtain ⟨k, hk⟩ := exists_nat_gt (max ((max X₀ 4 : ℕ) : ℝ) (1 / c))
  have hknat : max X₀ 4 < k := by exact_mod_cast (le_max_left _ _).trans_lt hk
  have hk4 : 4 ≤ k := by omega
  have hck : 1 < (k : ℝ) * c :=
    (div_lt_iff₀ hc).1 ((le_max_right _ _).trans_lt hk)
  have hsquare : k ^ 2 ≤ 2 ^ k * target :=
    (square_le_two_pow hk4).trans (Nat.le_mul_of_pos_right _ hpos)
  have hroom : X₀ ≤ 2 ^ k * target := by
    have hk_le_square : k ≤ k ^ 2 := by nlinarith
    omega
  have hcount := hX₀ (2 ^ k * target) hroom
  rw [three_divisible_count_at_dyadic_cutoff hpos hthree k] at hcount
  have hsquare_real : (k : ℝ) ^ 2 ≤ ((2 ^ k * target : ℕ) : ℝ) := by
    exact_mod_cast hsquare
  have hpaid := mul_le_mul_of_nonneg_left hsquare_real hc.le
  have hkpos : (0 : ℝ) < k := by exact_mod_cast (show 0 < k by omega)
  nlinarith

end CollatzTargetDensity
\end{Verbatim}
\subsection{\texorpdfstring{\texttt{CollatzTargetDensity/Public.lean}}{CollatzTargetDensity/Public.lean}}
\begin{Verbatim}[fontsize=\footnotesize,breaklines=true,breakanywhere=true,numbers=left,numbersep=5pt]
import CollatzTargetDensity.AllTargets
import CollatzTargetDensity.ThreeDivisible

/-! Literal map and count interface for the all-target density extension. -/
namespace CollatzTargetDensity

def step (n : ℕ) : ℕ := if Even n then n / 2 else 3 * n + 1

def reachesTarget (target n : ℕ) : Prop :=
  0 < n ∧ ∃ k : ℕ, (step^[k]) n = target

open Classical in
noncomputable def countBelow (target X : ℕ) : ℕ :=
  Nat.count (reachesTarget target) X

/-- The count uses positive starting integers strictly below X and the
    ordinary Collatz map. The lower-density constant may depend on target. -/
theorem positive_density (target : ℕ) (_hpos : 0 < target) (hthree : target % 3 ≠ 0) :
    ∃ c : ℝ, 0 < c ∧ ∃ X₀ : ℕ, ∀ X : ℕ, X₀ ≤ X →
      c * (X : ℝ) ≤ (countBelow target X : ℝ) := by
  exact positive_density_all_targets hthree

theorem positive_density_iff (target : ℕ) (hpos : 0 < target) :
    (∃ c : ℝ, 0 < c ∧ ∃ X₀ : ℕ, ∀ X : ℕ, X₀ ≤ X →
      c * (X : ℝ) ≤ (countBelow target X : ℝ)) ↔ target % 3 ≠ 0 := by
  constructor
  · intro h hthree
    exact no_positive_lower_density_three_divisible hpos hthree h
  · exact positive_density target hpos

end CollatzTargetDensity
\end{Verbatim}
\subsection{\texorpdfstring{\texttt{CollatzTargetDensity/OrbitInvariance.lean}}{CollatzTargetDensity/OrbitInvariance.lean}}
\begin{Verbatim}[fontsize=\footnotesize,breaklines=true,breakanywhere=true,numbers=left,numbersep=5pt]
import CollatzTargetDensity.RawBridge

/-! Elementary orbit facts for interpreting the predecessor density result. -/
namespace CollatzTargetDensity
open Erdos1135

theorem reaches_one_forward {n a : ℕ} (hna : Reaches n a) (hn : Reaches n 1) :
    Reaches a 1 := by
  obtain ⟨d, hd⟩ := hna
  obtain ⟨t, ht⟩ := hn
  have hcycle : (collatzStep^[3]) 1 = 1 := by decide
  have hperiod : (collatzStep^[3 * d]) 1 = 1 := by
    simpa only [Nat.mul_comm] using periodic_multiples hcycle d
  refine ⟨t + 2 * d, ?_⟩
  rw [← hd, ← Function.iterate_add_apply,
    show t + 2 * d + d = 3 * d + t by omega,
    Function.iterate_add_apply, ht]
  exact hperiod

theorem reaches_three_unit (n : ℕ) (hn : 0 < n) :
    ∃ a : ℕ, a % 3 ≠ 0 ∧ Reaches n a := by
  induction n using Nat.strong_induction_on with
  | h n ih =>
    by_cases hthree : n % 3 ≠ 0
    · exact ⟨n, hthree, reaches_refl n⟩
    · by_cases he : Even n
      · have hmod := Nat.even_iff.mp he
        have hpos : 0 < n / 2 := by omega
        have hlt : n / 2 < n := Nat.div_lt_self hn (by decide)
        obtain ⟨a, ha, hreach⟩ := ih (n / 2) hlt hpos
        refine ⟨a, ha, reaches_of_reaches_step ?_⟩
        simpa only [collatzStep_eq_div_two_of_even he] using hreach
      · exact ⟨3 * n + 1, by omega,
          reaches_of_step_eq (collatzStep_eq_three_mul_add_one_of_not_even he)⟩

def rawBad : Set ℕ := {n | 0 < n ∧ ¬Reaches n 1}

theorem rawReaches_subset_bad {a : ℕ} (ha : ¬Reaches a 1) :
    rawReaches a ⊆ rawBad := by
  intro n hn
  exact ⟨hn.1, fun h => ha (reaches_one_forward hn.2 h)⟩

theorem bad_three_unit_of_bad {n : ℕ} (hn : n ∈ rawBad) :
    ∃ a : ℕ, a % 3 ≠ 0 ∧ ¬Reaches a 1 := by
  obtain ⟨a, ha, hreach⟩ := reaches_three_unit n hn.1
  exact ⟨a, ha, fun h => hn.2 (hreach.trans h)⟩

end CollatzTargetDensity
\end{Verbatim}
\subsection{\texorpdfstring{\texttt{CollatzTargetDensity/CounterexampleDichotomy.lean}}{CollatzTargetDensity/CounterexampleDichotomy.lean}}
\begin{Verbatim}[fontsize=\footnotesize,breaklines=true,breakanywhere=true,numbers=left,numbersep=5pt]
import CollatzTargetDensity.AllTargets
import CollatzTargetDensity.OrbitInvariance

/-!
# Consequence of the all-target density theorem

It does not establish that the sparse-counterexample criterion holds.
-/
namespace CollatzTargetDensity
open Erdos1135
noncomputable section

/-- One counterexample would force a positive lower natural density of
    counterexamples. The density constant is not asserted to be uniform. -/
theorem counterexamples_positive_density (hbad : rawBad.Nonempty) :
    ∃ c : ℝ, 0 < c ∧ ∃ X₀ : ℕ, ∀ X : ℕ, X₀ ≤ X →
      c * (X : ℝ) ≤ (Terras.natCount rawBad X : ℝ) := by
  obtain ⟨n, hn⟩ := hbad
  obtain ⟨a, ha, hnot⟩ := bad_three_unit_of_bad hn
  obtain ⟨c, hc, X₀, hX₀⟩ := positive_density_all_targets ha
  refine ⟨c, hc, X₀, ?_⟩
  intro X hX
  have hmono := Terras.natCount_mono (rawReaches_subset_bad hnot) X
  exact (hX₀ X hX).trans (by exact_mod_cast hmono)

/-- The counterexample proportion becomes arbitrarily small along arbitrarily
    large cutoffs. This is a condition, not an established fact. -/
def SparseBadSubsequence : Prop :=
  ∀ ε : ℝ, 0 < ε → ∀ X₀ : ℕ, ∃ X : ℕ, X₀ ≤ X ∧ 0 < X ∧
    (Terras.natCount rawBad X : ℝ) < ε * (X : ℝ)

theorem collatz_iff_sparse_bad_subsequence :
    (∀ n : ℕ, 0 < n → Reaches n 1) ↔ SparseBadSubsequence := by
  constructor
  · intro hall ε hε X₀
    have hempty : rawBad = ∅ := by
      ext n
      constructor
      · intro hn
        exact False.elim (hn.2 (hall n hn.1))
      · intro hn
        exact False.elim hn
    let X := max X₀ 1
    have hXpos : 0 < X := by dsimp [X]; omega
    have hcount : Terras.natCount rawBad X = 0 := by
      rw [hempty]
      simp [Terras.natCount]
    refine ⟨X, le_max_left _ _, hXpos, ?_⟩
    rw [hcount, Nat.cast_zero]
    exact mul_pos hε (by exact_mod_cast hXpos)
  · intro hsparse n hn
    by_contra hnot
    obtain ⟨c, hc, X₀, hX₀⟩ := counterexamples_positive_density ⟨n, hn, hnot⟩
    obtain ⟨X, hX, _, hsmall⟩ := hsparse c hc X₀
    exact (not_lt_of_ge (hX₀ X hX)) hsmall

end
end CollatzTargetDensity
\end{Verbatim}

\end{document}
